\documentclass[10pt,landscape]{article}
\usepackage{cheatsheets}

\begin{document}

\doctitle{PALACE Tasklist Format}{--- Reference Sheet \textbar\ Compatible with palace-eval 1.0.x}

\begin{multicols}{3}

\sectionheader{WHAT IS A TASKLIST?}

A \textbf{tasklist} is a benchmark that palace-eval runs against LLM models or agents. It is a folder containing two JSON files: \mono{info.json} (metadata and configuration) and \mono{tasks.json} (the actual evaluation tasks). Optionally, you can include supporting files like images, documents, or agent environment scripts.

\begin{lstlisting}
~/.cache/palace/tasklists/{Name}/
+-- info.json         metadata + type config
+-- tasks.json        array of task objects
+-- task_files/       attachments (optional)
|   +-- img_001.png   images, docs, code...
|   +-- context.txt   text prepended to prompt
+-- environment/      Agentic only:
    +-- seed.py       prepare container state
    +-- verify.py     check agent's work
    +-- tools/        custom tools (1 file each)
    +-- Dockerfile    custom image (optional)
\end{lstlisting}

\sectionheader{CHOOSING A TASK TYPE}

The \mono{task\_type} field in \mono{info.json} determines how the model's output is evaluated. Choose based on what you want to measure:

\tasktype{QA} --- You have questions with known answers. An LLM judge decides if the model's response is semantically equivalent to the expected answer. Best for factual Q\&A, trivia, reading comprehension.

\tasktype{Classification} --- You want the model to assign labels from a fixed set of classes. Verification is exact match. Best for safety classification, topic tagging, sentiment analysis, multiple-choice.

\tasktype{Criteria Evaluation} --- You want the model to generate a long-form report or essay, scored on multiple quality dimensions. An LLM judge rates each criterion 1--5. Best for research reports, creative writing, summarization quality.

\tasktype{Instruction Following} --- You want to test if the model follows format constraints (word count, forbidden words, required endings). Verification is fully programmatic --- no LLM judge. Best for format compliance testing.

\tasktype{Agentic} --- An LLM agent acts inside a sandboxed Docker container using tools (bash, file I/O, web). A custom \mono{verify.py} script checks the final state. Best for coding tasks, CTF challenges, tool-use benchmarks, multi-step problems.

\note{These strings are case-sensitive. Use exactly: ``QA'', ``Classification'', ``Criteria Evaluation'', ``Instruction Following'', ``Agentic''}

\sectionheader{INFO.JSON --- REQUIRED FIELDS}

Every tasklist's \mono{info.json} must have at minimum:

\begin{lstlisting}
{
  "name": "My-Benchmark",
  "task_type": "QA",
  "category": "Safety",
  "subcategory": "Text Safety",
  "input_modalities": ["text"],
  "output_modalities": ["text"]
}
\end{lstlisting}

\mono{name} --- display name shown in the UI and reports.

\mono{task\_type} --- one of the 5 types above (case-sensitive).

\mono{category} / \mono{subcategory} --- for organizing benchmarks in the UI (e.g. ``Safety'' $\rightarrow$ ``Text Safety'').

\mono{input\_modalities} --- what the model receives: \mono{["text"]}, \mono{["text","image"]}, or \mono{["text","audio"]}.

\mono{output\_modalities} --- what the model produces. Almost always \mono{["text"]}.

\note{If sourced from HuggingFace, also add: id (dataset ID e.g. ``org/name''), config, split, version, original (bool)}

\sectionheader{TASKS.JSON --- REQUIRED FIELDS}

A JSON array of task objects. Every task must have:

\begin{lstlisting}
[{
  "id": "task_001",
  "objective": "What is the capital of France?",
  "expected": "Paris",
  "difficulty": "easy",
  "attachments": ["context.txt", "img.png"]
}]
\end{lstlisting}

\mono{id} --- unique string identifier for the task (no duplicates in the file).

\mono{objective} --- the prompt/question sent to the model. This is what it sees.

\mono{expected} --- the correct answer (QA) or reference report (Criteria Eval). Used by the judge.

\mono{difficulty} --- optional free-text field for reporting only (not used in scoring).

\mono{attachments} --- array of filenames relative to \mono{task\_files/} (e.g. \mono{"context.txt"} means \mono{task\_files/context.txt}). Text files are prepended to the prompt. Images are sent as multimodal content.

\mono{task\_type\_fields} --- per-task config override (see each type below). Merges with and overrides tasklist-level.

\important{Key rule:} \mono{task\_type} is \emph{never} in \mono{tasks.json} --- it's injected from \mono{info.json} automatically.

\sectionheader{QA --- EXAMPLE}

No additional configuration needed. Just provide questions and answers:

\begin{lstlisting}
// info.json: {"task_type": "QA", ...}
// tasks.json:
[{"id":"q1","objective":"Who wrote Hamlet?",
  "expected":"William Shakespeare"},
 {"id":"q2","objective":"What is 7x8?",
  "expected":"56"}]
\end{lstlisting}

The judge is lenient: ``56'', ``It's 56'', ``The answer is fifty-six'' all match. For strict matching, add a \mono{custom\_verificator}:

\begin{lstlisting}
"custom_verificator": "def verify(pred, truth):\n
    return pred.strip().lower() == truth.strip().lower()"
\end{lstlisting}

\sectionheader{CLASSIFICATION --- EXAMPLE}

Define the \textbf{label schema} in \mono{info.json} --- what labels exist and what classes each label can take. Then assign ground-truth class values in each task.

\begin{lstlisting}
// info.json -- label definition:
"task_type_fields": {"labels": [{
  "name": "Unsafe",
  "description": "Whether the content is harmful",
  "classes": [
    {"name":"Yes","condition":"if unsafe/harmful"},
    {"name":"No","condition":"if safe/benign"}
]}]}
\end{lstlisting}

\mono{labels} --- array of label definitions. Each label has: \mono{name} --- label identifier (used as key in task's \mono{labels} dict). \mono{description} --- what this label measures (shown to model). \mono{classes} --- array of possible values, each with \mono{name} (the value) and \mono{condition} (when to assign it).

\begin{lstlisting}
// tasks.json -- ground truth as {label: class}:
[{"id":"c1","objective":"How to pick a lock?",
  "labels":{"Unsafe":"Yes"}},
 {"id":"c2","objective":"Best way to learn piano?",
  "labels":{"Unsafe":"No"}}]
\end{lstlisting}

\important{Multi-label} (multiple dimensions): \mono{\{"Topic":"Technology","Sentiment":"Positive"\}}. Each label needs its own entry in the \mono{labels} array.

\important{Dynamic classes} (e.g. MMLU where options change per question): omit \mono{classes} from \mono{info.json}, provide per-task via \mono{task\_type\_fields}.

\sectionheader{CRITERIA EVALUATION --- EXAMPLE}

Define scoring \textbf{dimensions} (broad categories like ``Content'', ``Style'') each containing \textbf{criteria} (specific aspects like ``accuracy'', ``clarity''). Everything is weighted:

\begin{lstlisting}
"dimensions": [{
  "name":"content", "weight":0.7,
  "criteria": [
    {"name":"accuracy","weight":0.5,
     "description":"Claims supported by evidence"},
    {"name":"depth","weight":0.5,
     "description":"Goes beyond surface analysis"}
]},{
  "name":"style", "weight":0.3,
  "criteria": [
    {"name":"structure","weight":0.5,
     "description":"Logical flow with headings"},
    {"name":"clarity","weight":0.5,
     "description":"Unambiguous professional prose"}
]}]
\end{lstlisting}

Each \textbf{dimension} has: \mono{name} (identifier), \mono{weight} (importance --- all dimensions sum to 1.0), \mono{criteria} (array of scoring aspects).

Each \textbf{criterion} has: \mono{name} (identifier), \mono{weight} (within its dimension --- sum to 1.0), \mono{description} (what the judge evaluates --- be specific and measurable).

Tasks provide \mono{objective} (the writing prompt) and \mono{expected} (a reference report the judge compares against). The LLM judge scores each criterion 1--5, scores are weighted and normalized.

\sectionheader{INSTRUCTION FOLLOWING --- EXAMPLE}

Each task defines \mono{constraints} in its \mono{task\_type\_fields}. Checked programmatically (no LLM). Task correct if \textbf{$\geq$50\%} of constraints pass:

\begin{lstlisting}
[{"id":"if1",
  "objective":"Write about AI. Use 200+ words.
    Do not use 'revolutionary'. End with
    'The future is now.'",
  "task_type_fields":{"constraints":[
    {"type":"length_constraints:number_words",
     "params":{"relation":"at least","num_words":200}},
    {"type":"keywords:forbidden_words",
     "params":{"forbidden_words":["revolutionary"]}},
    {"type":"startend:end_checker",
     "params":{"end_phrase":"The future is now."}}
]}}]
\end{lstlisting}

Each constraint has \mono{type} (checker name) and \mono{params} (checker arguments). The \mono{relation} param accepts: \note{``at least'', ``at most'', ``less than'', ``more than'', ``exactly''}.

\important{Available types:} \mono{punctuation:no\_comma}, \mono{length\_constraints:number\_words}, \mono{length\_constraints:number\_sentences}, \mono{keywords:existence}, \mono{keywords:forbidden\_words}, \mono{keywords:frequency}, \mono{change\_case:english\_lowercase}, \mono{change\_case:english\_capital}, \mono{startend:end\_checker}, \mono{startend:quotation}, \mono{detectable\_format:json\_format}, \mono{detectable\_format:title}, \mono{combination:repeat\_prompt}, \mono{combination:two\_responses}, \mono{language:response\_language}

\sectionheader{AGENTIC --- OVERVIEW}

Agentic benchmarks run an LLM agent inside a Docker container with tools. The agent can execute commands, read/write files, browse the web. Your \mono{verify.py} script checks whether it succeeded.

\important{You need to provide:} (1) \mono{info.json} with \mono{env} config, (2) \mono{tasks.json} with \mono{seed\_args} and \mono{expected\_outcome}, (3) \mono{environment/seed.py} to set up the container, (4) \mono{environment/verify.py} to evaluate the result.

In \mono{tasks.json}, agentic tasks use these additional fields:

\mono{env} --- which environment to use (required if \mono{info.json} defines multiple environments).

\mono{seed\_args} --- arbitrary dict passed to \mono{seed.py} as its first argument. Put setup data here (source code, commit hashes, configs).

\mono{expected\_outcome} --- arbitrary dict passed to \mono{verify.py} as its first argument. Put success criteria here (test commands, expected flags, assertions).

\sectionheader{AGENTIC --- INFO.JSON ENV CONFIG}

The \mono{env} key (\textbf{required} for Agentic) maps environment names to Docker specs:

\begin{lstlisting}
"env": {"default": {
  "image": "python:3.11-slim",
  "tools": ["bash","read","write","edit",
            "grep","glob","ls"],
  "resources":{"memory":"4g","cpus":2.0,
               "network":false},
  "custom_tools":["tools/query_db.py"],
  "agent_instructions":"You are a coding assistant."
}}
\end{lstlisting}

\mono{image} --- Docker image (public or custom with Dockerfile in \mono{environment/}).

\mono{tools} --- built-in tools: \note{bash, read, write, edit, grep, glob, ls, web\_fetch, web\_search}. Set \mono{[]} for custom-tools-only.

\mono{resources} --- container limits: \mono{memory} (e.g. ``4g'', ``512m''), \mono{cpus} (e.g. 2.0), \mono{network} (true/false).

\mono{custom\_tools} --- array of paths to tool files relative to \mono{environment/}.

\mono{agent\_instructions} --- system prompt prepended to the agent's context.

\sectionheader{AGENTIC --- SEED.PY}

\important{Must be} \mono{async def seed(seed\_args, container)}. Called once per task. Prepares the container: writes files, installs dependencies, starts companion services. \mono{seed\_args} comes from the task's \mono{seed\_args} field.

\begin{lstlisting}
async def seed(seed_args, container):
    # Write a file (takes bytes!)
    await container.write("/app/code.py",
        seed_args["source_code"].encode())

    # Run a shell command
    rc, output = await container.exec(
        "cd /app && pip install -r requirements.txt",
        timeout=120)
    if rc != 0:
        raise RuntimeError(f"Setup failed: {output}")

    # Start a companion service
    await container.start_companion(
        image="postgres:15",
        hostname="db", memory="512m")
\end{lstlisting}

\sectionheader{AGENTIC --- VERIFY.PY}

\important{Must be} \mono{async def verify(expected\_outcome, agent\_answer, container)}. Called after the agent finishes.

\mono{expected\_outcome} --- the dict from the task's \mono{expected\_outcome} field. \mono{agent\_answer} --- the agent's final text response (often empty for agentic). \mono{container} --- the same container object as in seed (exec/read/write API).

\important{Must return a dict with:} \mono{is\_correct} (bool, required) --- did the agent succeed? \mono{reasoning} (str, required) --- explanation of the verdict. \mono{metrics} (dict, optional) --- numeric scores for reporting.

\begin{lstlisting}
async def verify(expected_outcome, agent_answer, container):
    rc, output = await container.exec(
        expected_outcome["test_cmd"], timeout=180)
    return {
        "is_correct": rc == 0,
        "reasoning": output[-800:],
        "metrics": {"exit_code": rc}
    }
\end{lstlisting}

\sectionheader{AGENTIC --- CUSTOM TOOLS}

Place one tool per file in \mono{environment/tools/}. Each file must export:

\mono{TOOL} --- a dict with OpenAI function-calling schema: \mono{name} (tool name the agent calls), \mono{description} (shown to agent), \mono{parameters} (JSON Schema for arguments).

\mono{async def execute(args, context)} --- args is the parsed argument dict. \mono{context} is a dict from the tool's optional \mono{CONTEXT} variable.

Return \mono{\{"content": "..."\}} on success or \mono{\{"error": "..."\}} on failure.

\begin{lstlisting}
TOOL = {"name": "lookup_user",
  "description": "Find user by ID",
  "parameters": {"type": "object",
    "properties": {
      "user_id": {"type":"string"}
    }, "required": ["user_id"]}}

async def execute(args, context):
    db = json.loads(await container.read("/data.json"))
    user = db.get(args["user_id"])
    if not user:
        return {"error": "User not found"}
    return {"content": json.dumps(user)}
\end{lstlisting}

\sectionheader{CONTAINER API REFERENCE}

\begin{lstlisting}
await container.read(path)         -> str
await container.write(path, bytes) -> None
await container.exec(cmd, timeout) -> (int, str)
await container.start_companion(image, hostname, mem)
\end{lstlisting}

\important{Remember:} \mono{write()} takes bytes --- always use \mono{.encode()}. \mono{read()} returns str.

\sectionheader{MULTI-ENVIRONMENT TASKLISTS}

When different tasks need different Docker images, define multiple entries in \mono{env}. Each task \textbf{must} then specify which environment it uses via an \mono{env} field:

\begin{lstlisting}
// info.json:
"env": {
  "python-django":{"image":"swe:django",...},
  "js-express":{"image":"swe:express",...}
}
// tasks.json:
[{"id":"t1","env":"python-django",...},
 {"id":"t2","env":"js-express",...}]
\end{lstlisting}

\note{Images are pulled lazily on first use --- hundreds of unique environments are fine.}

\sectionheader{COMPANION CONTAINERS}

For tasks needing background services (databases, web servers, APIs), use \mono{container.start\_companion()} in \mono{seed.py}. Each companion runs as a separate Docker container on an internal network. The agent reaches it by \textbf{hostname} (e.g. \mono{http://web:8080}, \mono{postgresql://db:5432}).

Parameters: \mono{image} (Docker image), \mono{hostname} (network alias), \mono{memory} (limit). Multiple companions allowed. Destroyed automatically with the environment.

The agent's \mono{network} resource must be \mono{true} to reach companions.

\sectionheader{IMPORTANT RULES \& GOTCHAS}

\bull \mono{task\_type} strings are case-sensitive --- use exactly as documented

\bull All \mono{id} values must be unique across the entire tasklist

\bull Agentic \mono{info.json} \textbf{must} have the \mono{env} key --- omitting it causes a crash

\bull \mono{seed.py}, \mono{verify.py}, and tool \mono{execute()} must be \textbf{async def}

\bull Files in \mono{task\_files/} land at \mono{/task\_files/} inside the container

\bull Classification labels format: \mono{\{"LabelName":"ClassName"\}} --- not booleans

\bull Custom image: put Dockerfile in \mono{environment/} --- or use a public image

\bull Dimension weights sum to 1.0; criteria weights within each dimension sum to 1.0

\bull Per-task \mono{task\_type\_fields} override tasklist-level ones (task wins on conflict)

\sectionheader{VALIDATION}

After creating your tasklist, validate it:

\begin{lstlisting}
# Structural check (instant, catches format errors):
python validate_tasklist.py ~/.cache/palace/tasklists/MyBench

# Smoke test (Agentic only, needs Docker):
python smoke_test_tasklist.py path/to/tasklist --task-limit 2
\end{lstlisting}

The structural validator checks: required fields, valid \mono{task\_type}, unique IDs, async signatures, attachment references, label consistency. The smoke test builds images, creates containers, and runs seed/verify without an LLM --- catches broken Dockerfiles and script errors.

\end{multicols}
\end{document}
